\makeabstract
    {
    Foliar Fog Water Uptake in Tall Conifers
    }
    {
     Wylee Drew\textsuperscript{1}, Cooper Merrill\textsuperscript{1}, Austin McGuire\textsuperscript{1}, Katie Nelson\textsuperscript{1}, Miranda Yang\textsuperscript{1}, Rebecca Hewitt\textsuperscript{1}, Lucy Nelson\textsuperscript{2}
    }
    {
    \textsuperscript{1} Department of Environmental Studies, Amherst College, Amherst, MA\\
\textsuperscript{2} Department of Forestry, Fire \& Rangeland Management, Cal Poly Humboldt, Arcata, CA
    }
    {
    Tall conifers have difficulty transporting water from the roots to the canopy due to their immense height. To compensate, they can take up fog water directly through their leaves, yet the mechanism of this remains unclear. It has been suggested that foliar fungal endophytes (FFEs) play a role, with fungal hyphae that protrude through leaf stomata and wick moisture into the leaf interior.
    }
    {
    It has been suggested that foliar fungal endophytes (FFEs) play a role, with fungal hyphae that protrude through leaf stomata and wick moisture into the leaf interior. To investigate this, we collected foliage samples from three conifer species{\textendash}Coastal Redwood (Sequoia sempervirens), Sitka Spruce (Picea sitchensis), and Douglas Fir (Pseudotsuga menziesii){\textendash}from the Redwood Experimental Forest near Klamath, California. Foliage samples were collected from old and young age groups at lower, middle, and upper canopy positions from three replicate trees per species. For each sample we measured leaf chemistry (nutrients and isotopes), water uptake capacity in a fog chamber, and morphology, i.e.,  stomatal density and wax-covered fraction of the abaxial (bottom) surface. To locate FFEs, we   stained samples with WGA-AF594 and imaged them using a confocal microscope, and we extracted DNA from ground foliage d to determine the taxonomy of fungi present. We are currently quantifying FFE abundance using an assay that measures the fluorescence of WGA-AF594 in ground leaf samples. 
    }
    {
     Consistent with prior work, water uptake capacity was higher in old foliage than young; we also found higher 151N in leaves with greater fog water uptake, consistent with uptake of 15N-enriched, marine-derived nitrogen. Stomatal density and wax -covered fraction were not correlated with fog water uptake, suggesting that fungi, rather than leaf traits, may drive uptake.
    }
    {
    Ongoing fungal quantification will test this possibility, and clarifying how these tall forests capture fog water is key to understanding their resilience to drought.
    }

    \newpage
    